% Pursuing Stacks (cotes 134-2 to 134-8): archive cote 134-6, site page 67 (PDF page 68), typescript page 385.
% typescript. Draft compared: ps5.tex:4756-4839.
% AI-transcribed by Claude Opus 5.5 (claude-opus-5-5, adaptive thinking, effort high), Message Batches API, batch ps3 (msgbatch_01XdFmuVVBAcgc7EFxdtCqsj), 2026-11-12; not yet checked by eye.
\scanpage{134-6/67}
\mspage{385}

For the sake of brevity, we'll say that a \uline{U}-category satisfying the
assumptions of prop.1 is a \uline{pseudo}-\uline{topos} (as these conditions are satisfied
for any topos). We get at once the

\uline{Corollary} 2. Let \hand{$\mathcal{A}$},\hand{$\mathcal{B}$} be two pseudo-topoi. Then a functor from
\hand{$\mathcal{A}$} to \hand{$\mathcal{B}$}$^{o}$ (resp. from \hand{$\mathcal{B}$}$^{o}$ to \hand{$\mathcal{A}$} ) has a right adjoint (resp. a left
adjoint) iff it commutes to small direct limits (resp. to small inverse
limits). Thus taking right and left adjoints we get two equivalences of
categories, quasi-inverse to each other
\[ (4) \qquad \underline{Hom}_{!}(\hand{\mathcal{A}},\hand{\mathcal{B}}^{o})\hand{{}^{o}}\hand{\leftrightarrows}\underline{Hom}^{!}(\hand{\mathcal{B}}^{o},\hand{\mathcal{A}})\deleted{\hand{{}^{o}}} \quad , \]
and the two members of (4) are canonically equivalent to the \deleted{\unsure{dual of the}} category
\[ (5) \qquad \underline{Hom}^{\hand{!!}}_{\deleted{\illegible{1}}}(\hand{\mathcal{A}^{o}},\hand{\mathcal{B}^{o}};\text{ (Sets)}) \]
of functors
\[ \hand{\mathcal{A}^{o}}\text{x}\hand{\mathcal{B}^{o}\to}\text{(Sets)} \]
which commute with small \xout{direct}\inserted[typed]{inverse} limits with respect to either variable
(the other being fixed),\inserted{(5) being} viewed as a full subcategory of \uline{Hom}(\hand{$\mathcal{A}^{o}$}x\hand{$\mathcal{B}^{o}$},(Sets\hand{)})

\xout{Thexbifunctorxxxx$\overset{o}{\text{x}}$Bxxxxx(Sets)xxassociatedxtoxaxfunctor}
\uline{Remarks}. \hand{1)}The first hand side of (4)\deleted{$_{o}$} is tautologically isomorphic to
the category Hom$^{!}$(\hand{$\mathcal{A}$}$^{o}$,\deleted{\illegible{1}}\hand{$\mathcal{B}$}) (as for any two\inserted[typed]{U-}categories we have the
\xout{is of course} tautological \deleted{equivalence} \hand{isomorphism}
\[ (6) \qquad \deleted{(x,y)xx} \quad (\underline{Hom}_{!}(\hand{\mathcal{P}},\hand{\mathcal{Q}}))^{o}\hand{\simeq}\underline{Hom}^{!}(\hand{\mathcal{P}^{o}},\hand{\mathcal{Q}^{o}}) \quad ) \quad , \]
thus the equivalence (4) can be seen more symetrically as an equivalence
\[ (4') \qquad \underline{Hom}^{!}(\hand{\mathcal{A}}^{o},\hand{\mathcal{B}})\hand{\unsure{\cong|\simeq}}\underline{Hom}^{!}(\hand{\mathcal{B}}^{o},\hand{\mathcal{A}}) \quad , \]
both categories being equivalent to (5) by using the equivalence (3)
for the second, and the corresponding equivalence for \hand{$\mathcal{B}$} for the first,
plus the tautological isomorphism
\[ (7) \qquad \underline{Hom}^{!}(\hand{\mathcal{P}},\underline{Hom}^{!}(\hand{\mathcal{Q}},\hand{\mathcal{M}})) \qquad \underline{Hom}^{!!}(\hand{\mathcal{P},\mathcal{Q};\mathcal{M}}) \quad , \]
for any three \uline{U}-categories \hand{$\mathcal{P}$},\hand{$\mathcal{Q}$},\hand{$\mathcal{M}$}.

\hand{2)} When \hand{$\mathcal{A}$} is a topos, \hand{$\mathcal{B}$} any \uline{U}-category stable under\inserted[typed]{small} inverse limits,
then we may interprete the category \uline{Hom}$^{!}$(\hand{$\mathcal{A}$}$^{o}$,\hand{$\mathcal{B}$}) as the category of
\hand{$\mathcal{B}$}-\uline{valued} \uline{sheaves} on the topos (defined by)\hand{$\mathcal{A}$}. When \hand{$\mathcal{A}$} and \hand{$\mathcal{B}$} are both
topoi, then the equivalence (4) states that \xout{shxx} \hand{$\mathcal{B}$}-valued sheaves on
\hand{$\mathcal{A}$} can be identified with \hand{$\mathcal{A}$}-valued sheaves on \hand{$\mathcal{B}$}, and both may be
identified with set-valued ``bi-sheaves'' on \hand{$\mathcal{A}$} x \hand{$\mathcal{B}$}. In case the topoi
\hand{$\mathcal{A}$},\hand{$\mathcal{B}$} are defined respectively by\inserted{\underline{U}-}sites A,B (not necessarily small ones),

\margin{67}
